\section{Posizionamento del contributo}
\label{sec:posizionamento}

Le tre famiglie di soluzioni analizzate in questo capitolo ---
scheduling energy-aware, approximate computing e carbon-aware computing
--- operano su dimensioni distinte del problema: la scelta del contesto
di esecuzione, la modulazione della logica applicativa e la sensibilità
al segnale ambientale. Tuttavia, nessuno degli approcci esaminati le
integra in un unico meccanismo decisionale.

In particolare, nessuno degli approcci esaminati affronta il problema
di selezionare, al momento dell'invocazione, la variante di una
funzione che offra il miglior compromesso tra consumo energetico e
accuratezza del risultato in funzione dell'impatto ambientale corrente
dell'esecuzione. Manca un meccanismo che utilizzi la carbon intensity
non come criterio di \emph{placement} geografico o temporale, ma come
segnale per modulare il punto operativo sul trade-off
energia-accuratezza della funzione stessa --- spostando la selezione
verso varianti a minor consumo quando la rete è alimentata da fonti ad
alta intensità di carbonio, e verso varianti a maggiore accuratezza
quando le condizioni della rete lo consentono.

A ciò si aggiunge una seconda lacuna, trasversale agli approcci
esaminati: in ambienti FaaS, il consumo energetico di una funzione non
è una grandezza stabile nel tempo, ma varia in funzione del carico
concorrente sul nodo, dello stato del container e delle caratteristiche
dei dati di input. Fondare le decisioni di selezione su profili
energetici determinati offline significa operare su stime
potenzialmente disallineate rispetto alle condizioni operative correnti,
con conseguenze dirette sull'efficacia del meccanismo decisionale.

Il presente lavoro si propone di colmare questo spazio, integrando le
tre dimensioni identificate in un unico sistema di scheduling che
operi a livello di singola invocazione, selezionando tra varianti della
stessa funzione sulla base di stime energetiche aggiornate
continuamente e di un segnale di carbon intensity contestualizzato
rispetto alla zona geografica di esecuzione.